% ============================================================
%  Semaine 3 — HTML5 : formulaires & accessibilité  [Module B]
%  PRÉSENTATION (Beamer) — cfpt·i | Anne Terrier
%  Aligné sur PlanCours_Pass_Web_final et sur S03_polycop
%
%  Images attendues dans ./images/ : (aucune obligatoire)
%  Une image absente est remplacée par un cadre de réservation.
% ============================================================
\documentclass[aspectratio=169,11pt]{beamer}

\usepackage[T1]{fontenc}
\usepackage[utf8]{inputenc}
\usepackage[french]{babel}
\usepackage[scaled=0.92]{helvet}
\renewcommand{\familydefault}{\sfdefault}

\usepackage{graphicx}
\usepackage{tikz}
\usetikzlibrary{positioning, arrows.meta, decorations.pathreplacing}
\usepackage{booktabs}
\usepackage{tabularx}
\usepackage{colortbl}
\usepackage{listings}

% ---------------------------------------------------------------
% Correction de l'alignement des en-tetes de tableau
%   1) \aboverulesep / \belowrulesep a 0 : supprime le lisere blanc
%      entre le filet booktabs et la bande coloree du \rowcolor.
%   2) \hdr : force l'alignement de la cellule d'en-tete. Sans cela,
%      une colonne p{}, m{} ou X place le texte trop bas dans la bande.
% ---------------------------------------------------------------
\setlength{\aboverulesep}{0pt}
\setlength{\belowrulesep}{0pt}
\newcommand{\hdr}[1]{\multicolumn{1}{l}{\color{white}\bfseries #1}}
\newcommand{\hdrc}[1]{\multicolumn{1}{c}{\color{white}\bfseries #1}}
\newcommand{\hdrr}[1]{\multicolumn{1}{r}{\color{white}\bfseries #1}}

% ---------------------------------------------------------------
% Palette cfpt·i
% ---------------------------------------------------------------
\definecolor{cfptYellow}{HTML}{F5C400}
\definecolor{cfptBlack}{HTML}{1A1A1A}
\definecolor{cfptDarkGrey}{HTML}{555555}
\definecolor{cfptGrey}{HTML}{F2F2F2}
\definecolor{accentBlue}{HTML}{1D4E89}
\definecolor{lightBlue}{HTML}{D6E4F0}
\definecolor{codeGrey}{HTML}{F5F5F5}
\definecolor{successGreen}{HTML}{1E7145}
\definecolor{warningOrange}{HTML}{C55A11}
\definecolor{alertRed}{HTML}{A32020}

% Alias (repris du polycopié)
\colorlet{uniInk}{cfptBlack}
\colorlet{uniNavy}{accentBlue}
\colorlet{uniTeal}{successGreen}
\colorlet{uniGold}{warningOrange}
\colorlet{uniGrey}{cfptDarkGrey}
\colorlet{uniPaper}{cfptGrey}
\colorlet{uniRule}{cfptDarkGrey!35}

% ---------------------------------------------------------------
% Image robuste : cadre de réservation si le fichier est absent
% ---------------------------------------------------------------
\newcommand{\imageslide}[2][0.80\textwidth]{%
  \IfFileExists{#2}%
    {\includegraphics[width=#1]{#2}}%
    {\setlength{\fboxrule}{0.6pt}%
     \textcolor{cfptDarkGrey!60}{\fbox{\parbox[c][3.6cm][c]{#1}{%
       \centering\small\color{cfptDarkGrey}%
       [ image à insérer ]\\[4pt]\texttt{\detokenize{#2}}}}}}%
}

% ---------------------------------------------------------------
% Thème Beamer cfpt·i
% ---------------------------------------------------------------
\usetheme{default}
\usecolortheme{default}
\setbeamercolor{palette primary}{bg=accentBlue, fg=white}
\setbeamercolor{structure}{fg=accentBlue}
\setbeamercolor{frametitle}{bg=accentBlue, fg=white}
\setbeamercolor{title}{fg=white}
\setbeamercolor{block title}{bg=accentBlue, fg=white}
\setbeamercolor{block body}{bg=lightBlue!30, fg=cfptBlack}
\setbeamercolor{block title alerted}{bg=cfptYellow!85!black, fg=cfptBlack}
\setbeamercolor{block body alerted}{bg=cfptYellow!20, fg=cfptBlack}
\setbeamercolor{block title example}{bg=successGreen, fg=white}
\setbeamercolor{block body example}{bg=successGreen!8, fg=cfptBlack}
\setbeamercolor{itemize item}{fg=cfptYellow!85!black}
\setbeamercolor{itemize subitem}{fg=cfptDarkGrey}
\setbeamercolor{enumerate item}{fg=accentBlue}
\setbeamercolor{normal text}{fg=cfptBlack}
\setbeamercolor{footline}{fg=white}

\setbeamertemplate{itemize item}{\textbullet}
\setbeamertemplate{itemize subitem}{--}
\setbeamertemplate{navigation symbols}{}
\setbeamertemplate{blocks}[rounded][shadow=false]

% Titre de frame : bandeau pleine largeur
\setbeamertemplate{frametitle}{%
  \nointerlineskip%
  \begin{beamercolorbox}[wd=\paperwidth, ht=2.6ex, dp=1.2ex, leftskip=0.3cm]{frametitle}%
    \usebeamerfont{frametitle}\bfseries\insertframetitle%
  \end{beamercolorbox}%
}

% Pied de page
\setbeamertemplate{footline}{%
  \leavevmode%
  \hbox{%
    \begin{beamercolorbox}[wd=.32\paperwidth, ht=2.4ex, dp=1ex, leftskip=0.3cm]{palette primary}%
      \usebeamerfont{author in head/foot}\tiny Anne Terrier%
    \end{beamercolorbox}%
    \begin{beamercolorbox}[wd=.44\paperwidth, ht=2.4ex, dp=1ex, center]{palette primary}%
      \usebeamerfont{title in head/foot}\tiny Semaine 3 --- HTML5 : formulaires \& accessibilité%
    \end{beamercolorbox}%
    \begin{beamercolorbox}[wd=.24\paperwidth, ht=2.4ex, dp=1ex, right, rightskip=0.3cm]{palette primary}%
      \usebeamerfont{date in head/foot}\tiny\insertframenumber\,/\,\inserttotalframenumber%
    \end{beamercolorbox}%
  }%
}

% Diapositive de séparation de partie
\newcommand{\partieslide}[1]{%
{%
\setbeamertemplate{footline}{}%
\begin{frame}[plain]
  \begin{tikzpicture}[remember picture, overlay]
    \fill[accentBlue] (current page.north west) rectangle (current page.south east);
    \fill[cfptYellow] ([yshift=-0.15cm]current page.center -| current page.west)
      rectangle ([yshift=-0.3cm]current page.center -| current page.east);
    \node[text=white, font=\Huge\bfseries] at ([yshift=0.9cm]current page.center) {#1};
  \end{tikzpicture}
\end{frame}
}}

% ---------------------------------------------------------------
% Listings
% ---------------------------------------------------------------
\lstdefinelanguage{HTML5}{
  keywords={<!DOCTYPE, html, head, body, header, nav, main, section, article,
    aside, footer, h1, h2, h3, h4, h5, h6, p, a, img, figure, figcaption,
    ul, ol, li, table, caption, thead, tbody, tr, th, td, form, fieldset,
    legend, input, textarea, select, option, optgroup, datalist, label,
    button, div, span, meta, title, link, script, strong, em, br, hr},
  sensitive=false,
  morecomment=[s]{<!--}{-->},
  morestring=[b]",
  keywordstyle=\color{accentBlue}\bfseries,
  stringstyle=\color{warningOrange},
  commentstyle=\color{cfptDarkGrey}\itshape,
}

\lstdefinelanguage{CSS}{
  morekeywords={background-color, border, border-radius, color, cursor,
    display, flex, font-family, font-size, grid, height, margin, outline,
    padding, text-align, width},
  morecomment=[s]{/*}{*/},
  morestring=[b]",
  sensitive=false,
}

\lstset{
  basicstyle=\ttfamily\scriptsize,
  backgroundcolor=\color{codeGrey},
  frame=single, framerule=0.4pt,
  rulecolor=\color{cfptDarkGrey!40},
  breaklines=true,
  keywordstyle=\color{accentBlue}\bfseries,
  commentstyle=\color{cfptDarkGrey}\itshape,
  stringstyle=\color{successGreen},
  showstringspaces=false,
  extendedchars=true,
  literate=%
    {é}{{\'e}}1 {è}{{\`e}}1 {ê}{{\^e}}1 {ë}{{\"e}}1
    {É}{{\'E}}1 {È}{{\`E}}1 {Ê}{{\^E}}1
    {à}{{\`a}}1 {â}{{\^a}}1 {À}{{\`A}}1
    {î}{{\^i}}1 {ï}{{\"i}}1 {ô}{{\^o}}1
    {û}{{\^u}}1 {ù}{{\`u}}1 {ü}{{\"u}}1
    {ç}{{\c c}}1 {Ç}{{\c C}}1 {œ}{{\oe}}1
}

% ---------------------------------------------------------------
% Métadonnées
% ---------------------------------------------------------------
\title{HTML5 --- Formulaires \& accessibilité}
\subtitle{Développement web --- Classe passerelle}
\author{Anne Terrier}
\institute{cfpt\textperiodcentered i --- École d'informatique}
\date{Semaine 3}

% ===============================================================
\begin{document}
% ===============================================================

% --- Diapo de titre ---
{
\setbeamertemplate{footline}{}
\begin{frame}[plain]
  \begin{tikzpicture}[remember picture, overlay]
    \fill[accentBlue] (current page.north west) rectangle ([yshift=-3.2cm]current page.north east);
    \node[anchor=north west, text=white, font=\huge\bfseries, text width=0.9\paperwidth]
      at ([shift={(0.8cm,-0.8cm)}]current page.north west)
      {HTML5 --- Formulaires \& accessibilité};
    \node[anchor=north west, text=lightBlue, font=\large]
      at ([shift={(0.85cm,-2.55cm)}]current page.north west)
      {Développement web --- Classe passerelle};

    % accent jaune
    \fill[cfptYellow] ([yshift=-3.2cm]current page.north west) rectangle ([yshift=-3.35cm]current page.north east);

    \node[anchor=south west, font=\large\bfseries, text=accentBlue]
      at ([shift={(0.8cm,1.8cm)}]current page.south west) {Semaine 3 \quad\textbar\quad Module B};
    \node[anchor=south west, font=\normalsize, text=cfptDarkGrey]
      at ([shift={(0.8cm,1.1cm)}]current page.south west)
      {Anne Terrier \quad\textbar\quad cfpt\textperiodcentered i --- École d'informatique};
  \end{tikzpicture}
\end{frame}
}

% --- Plan ---
\begin{frame}{Au programme aujourd'hui}
  \begin{columns}[T]
    \begin{column}{0.48\textwidth}
      \begin{enumerate}
        \item À quoi sert un formulaire
        \item Les champs de saisie
        \item Les attributs qui encadrent la saisie
      \end{enumerate}
    \end{column}
    \begin{column}{0.48\textwidth}
      \begin{enumerate}
        \setcounter{enumi}{3}
        \item La validation native et ses limites
        \item Structurer un formulaire
        \item Accessibilité
      \end{enumerate}
    \end{column}
  \end{columns}
  \vspace{0.4cm}
  \begin{block}{Objectif de la semaine}
    Écrire des formulaires complets, corrects et utilisables par tous --- y compris par une personne qui n'utilise ni souris ni écran.
  \end{block}
\end{frame}

\begin{frame}{Où en sommes-nous ?}
  \begin{columns}[T]
    \begin{column}{0.49\textwidth}
      \begin{block}{Semaine 2 --- afficher}
        Le MiniForum \textbf{présente} du contenu : titres, listes de sujets, messages. Tout est écrit en dur dans le HTML.
      \end{block}
    \end{column}
    \begin{column}{0.49\textwidth}
      \begin{exampleblock}{Semaine 3 --- saisir}
        La page s'ouvre dans l'autre sens : le visiteur va pouvoir \textbf{entrer} quelque chose.
      \end{exampleblock}
    \end{column}
  \end{columns}
  \vspace{0.35cm}
  \begin{alertblock}{Dernière séance de HTML}
    Les formulaires écrits aujourd'hui sont ceux du MiniForum \textbf{final} : leur structure ne changera plus jusqu'à la soutenance. Ils resteront inertes quelques semaines --- JavaScript les branchera en semaine 7, le serveur en semaine 15.
    \vspace{0.1cm}

    Le CSS commence la semaine prochaine.
  \end{alertblock}
\end{frame}

% ===============================================================
\section{À quoi sert un formulaire}
% ===============================================================
\partieslide{1. À quoi sert un formulaire}

\begin{frame}{Le point d'entrée des données}
  \begin{block}{Définition --- Formulaire}
    Un \textbf{formulaire} est un ensemble de champs de saisie, regroupés dans une balise \texttt{<form>}, permettant à un visiteur de transmettre des données à l'application.
  \end{block}
  \vspace{0.3cm}
  \begin{columns}[T]
    \begin{column}{0.49\textwidth}
      \begin{exampleblock}{Sans formulaire}
        Le site \textbf{diffuse}. Le contenu vient du développeur, et de lui seul.
      \end{exampleblock}
    \end{column}
    \begin{column}{0.49\textwidth}
      \begin{exampleblock}{Avec formulaire}
        Le site \textbf{reçoit}. Un forum, une boutique, un réseau social n'existent que par ce que les visiteurs y déposent.
      \end{exampleblock}
    \end{column}
  \end{columns}
\end{frame}

\begin{frame}{La vie d'une donnée saisie}
  \vspace{0.3cm}
  \begin{center}
  \begin{tikzpicture}[font=\scriptsize, scale=0.95, transform shape,
    et/.style={draw, rounded corners=2pt, minimum width=2.35cm,
               minimum height=1.4cm, align=center, inner sep=3pt},
    fl/.style={-{Stealth[length=5pt]}, thick, uniGrey!70}]

  \node[et, draw=uniNavy, fill=uniNavy!10] (a) at (0,0)
    {\textbf{Saisie}\\dans le\\formulaire};
  \node[et, draw=uniNavy, fill=uniNavy!10, right=0.42cm of a] (b)
    {\textbf{Contrôle}\\par le\\navigateur};
  \node[et, draw=uniGold, fill=uniGold!10, right=0.42cm of b] (c)
    {\textbf{Contrôle}\\en\\JavaScript};
  \node[et, draw=uniTeal, fill=uniTeal!10, right=0.42cm of c] (d)
    {\textbf{Envoi}\\au serveur\\et contrôle};
  \node[et, draw=uniTeal, fill=uniTeal!10, right=0.42cm of d] (e)
    {\textbf{Écriture}\\en base de\\données};

  \draw[fl] (a) -- (b);
  \draw[fl] (b) -- (c);
  \draw[fl] (c) -- (d);
  \draw[fl] (d) -- (e);

  \node[below=0.12cm of a, font=\scriptsize\bfseries, text=uniNavy] {semaine 3};
  \node[below=0.12cm of b, font=\scriptsize\bfseries, text=uniNavy] {semaine 3};
  \node[below=0.12cm of c, font=\scriptsize\bfseries, text=uniGold] {semaine 7};
  \node[below=0.12cm of d, font=\scriptsize\bfseries, text=uniTeal] {semaine 15};
  \node[below=0.12cm of e, font=\scriptsize\bfseries, text=uniTeal] {semaine 15};
  \end{tikzpicture}
  \end{center}
  \vspace{0.3cm}
  \begin{center}
    \small Aujourd'hui : les \textbf{deux premières cases}. Les trois autres expliquent pourquoi certains choix faits ce matin paraîtront inutiles jusqu'en décembre.
  \end{center}
\end{frame}

\begin{frame}[fragile]{La balise \texttt{<form>}}
\begin{lstlisting}[language=HTML5]
<form action="/topics" method="post">
  <!-- les champs de saisie viennent ici -->
  <button type="submit">Publier</button>
</form>
\end{lstlisting}
  \vspace{0.2cm}
  \begin{center}
  \small
  \begin{tabularx}{0.95\textwidth}{lX}
    \toprule
    \rowcolor{accentBlue}
    \hdr{Attribut} & \hdr{Rôle} \\
    \midrule
    \rowcolor{cfptGrey}
    \texttt{action} & L'adresse à laquelle les données sont envoyées \\
    \texttt{method} & La façon de les envoyer : \texttt{get} (dans l'URL) ou \texttt{post} (dans le corps de la requête) \\
    \bottomrule
  \end{tabularx}
  \end{center}
\end{frame}

\begin{frame}{\texttt{get} ou \texttt{post} ?}
  \begin{alertblock}{Ces deux attributs ne serviront qu'en semaine 15}
    Tant qu'aucun serveur n'écoute à l'adresse indiquée, \texttt{action} et \texttt{method} n'ont aucun effet visible : cliquer sur « Publier » recharge la page, et c'est tout.
  \end{alertblock}
  \vspace{0.2cm}
  \begin{columns}[T]
    \begin{column}{0.49\textwidth}
      \begin{block}{\texttt{get}}
        Les données passent \textbf{dans l'URL}.
        \vspace{0.1cm}

        Pratique pour une recherche que l'on veut mettre en favori.
        \vspace{0.1cm}

        \textcolor{alertRed}{\textbf{Inacceptable}} pour un mot de passe.
      \end{block}
    \end{column}
    \begin{column}{0.49\textwidth}
      \begin{exampleblock}{\texttt{post}}
        Les données passent dans le \textbf{corps de la requête}.
        \vspace{0.1cm}

        Invisibles dans la barre d'adresse.
        \vspace{0.1cm}

        Le choix par défaut dès qu'on crée ou modifie quelque chose.
      \end{exampleblock}
    \end{column}
  \end{columns}
\end{frame}

% ===============================================================
\section{Les champs de saisie}
% ===============================================================
\partieslide{2. Les champs de saisie}

\begin{frame}{\texttt{<input>} et son attribut \texttt{type}}
  Une seule balise, auto-fermante, dont le comportement change \textbf{entièrement} selon son \texttt{type}.
  \vspace{0.15cm}
  \begin{center}
  \scriptsize
  \begin{tabularx}{\textwidth}{lXl}
    \toprule
    \rowcolor{accentBlue}
    \hdr{\texttt{type}} & \hdr{Usage} & \hdr{Contrôle automatique} \\
    \midrule
    \rowcolor{cfptGrey}
    \texttt{text} & Texte libre sur une ligne & aucun \\
    \texttt{email} & Adresse e-mail & présence d'un \texttt{@} et d'un domaine \\
    \rowcolor{cfptGrey}
    \texttt{password} & Mot de passe & aucun ; masque les caractères \\
    \texttt{number} & Valeur numérique & chiffres, bornes \texttt{min} / \texttt{max} \\
    \rowcolor{cfptGrey}
    \texttt{date} & Date & calendrier, format garanti \\
    \texttt{url} & Adresse web & format d'URL \\
    \rowcolor{cfptGrey}
    \texttt{tel} & Numéro de téléphone & aucun (les formats varient trop) \\
    \texttt{checkbox} & Case à cocher & --- \\
    \rowcolor{cfptGrey}
    \texttt{radio} & Choix unique dans un groupe & --- \\
    \texttt{hidden} & Donnée transmise sans être affichée & --- \\
    \bottomrule
  \end{tabularx}
  \end{center}
\end{frame}

\begin{frame}{Le bon \texttt{type} améliore surtout le téléphone}
  \begin{exampleblock}{Un gain gratuit}
    Sur mobile, le clavier affiché s'adapte au \texttt{type} du champ :
    \begin{itemize}
      \item \texttt{number} $\rightarrow$ pavé numérique
      \item \texttt{email} $\rightarrow$ clavier avec la touche \texttt{@}
      \item \texttt{date} $\rightarrow$ sélecteur de calendrier
      \item \texttt{tel} $\rightarrow$ clavier de téléphone
    \end{itemize}
  \end{exampleblock}
  \vspace{0.3cm}
  \begin{center}
    \small Choisir le bon \texttt{type} est le moyen le \textbf{moins coûteux} d'améliorer un formulaire : un mot à changer, aucun code à écrire.
  \end{center}
\end{frame}

\begin{frame}[fragile]{Les autres champs}
\begin{lstlisting}[language=HTML5]
<!-- Texte long, sur plusieurs lignes -->
<textarea id="message" name="message" rows="6" maxlength="500"></textarea>

<!-- Liste deroulante -->
<select id="categorie" name="categorie">
  <option value="general">General</option>
  <option value="aide">Aide</option>
</select>

<button type="submit">Publier</button>
\end{lstlisting}
  \vspace{0.1cm}
  \begin{center}
  \scriptsize
  \begin{tabularx}{\textwidth}{lX}
    \toprule
    \rowcolor{accentBlue}
    \hdr{Balise} & \hdr{Rôle} \\
    \midrule
    \rowcolor{cfptGrey}
    \texttt{<textarea>} & \textbf{Jamais auto-fermante} : le contenu initial se place entre les deux balises \\
    \texttt{<select>} / \texttt{<option>} & C'est l'attribut \texttt{value} de l'\texttt{<option>} qui est transmis, pas le texte affiché \\
    \rowcolor{cfptGrey}
    \texttt{<button>} & \texttt{type="submit"} envoie le formulaire ; \texttt{type="button"} ne fait rien sans JavaScript \\
    \bottomrule
  \end{tabularx}
  \end{center}
\end{frame}

\begin{frame}{Le piège du \texttt{<button>} sans \texttt{type}}
  \begin{alertblock}{Le comportement par défaut}
    À l'intérieur d'un \texttt{<form>}, un \texttt{<button>} sans attribut \texttt{type} vaut \texttt{submit} : il \textbf{envoie le formulaire et recharge la page}.
  \end{alertblock}
  \vspace{0.3cm}
  \begin{block}{Ce que ça donne en semaine 7}
    C'est la cause d'une bonne moitié des « mon bouton ne marche pas » : le code JavaScript s'exécute bien, mais la page se recharge dans la foulée et efface tout.
    \vspace{0.15cm}

    \textbf{Précise toujours le \texttt{type}}, même quand il semble évident.
  \end{block}
\end{frame}

\begin{frame}[fragile]{Cases à cocher et boutons radio}
\begin{lstlisting}[language=HTML5]
<!-- Case a cocher : chaque case est independante -->
<input type="checkbox" id="news" name="news">
<label for="news">Recevoir la newsletter</label>

<!-- Boutons radio : meme name = meme groupe -->
<input type="radio" id="pub" name="visibilite" value="public" checked>
<label for="pub">Public</label>

<input type="radio" id="priv" name="visibilite" value="prive">
<label for="priv">Prive</label>
\end{lstlisting}
  \vspace{0.15cm}
  \begin{center}
    \small Une \textbf{case à cocher} autorise plusieurs sélections indépendantes ; un \textbf{bouton radio} impose un choix unique dans un groupe.
  \end{center}
\end{frame}

\begin{frame}{La règle des boutons radio}
  \begin{alertblock}{Même \texttt{name} pour tout le groupe}
    Pour qu'un groupe de boutons radio fonctionne --- un seul sélectionnable à la fois --- tous ses boutons doivent partager le \textbf{même \texttt{name}}.
  \end{alertblock}
  \vspace{0.2cm}
  \begin{center}
  \small
  \begin{tabularx}{0.95\textwidth}{lX}
    \toprule
    \rowcolor{accentBlue}
    \hdr{Attribut} & \hdr{Dans un groupe de radios} \\
    \midrule
    \rowcolor{cfptGrey}
    \texttt{name} & \textbf{identique} pour tous : c'est lui qui les relie \\
    \texttt{value} & \textbf{différente} pour chacun : c'est l'option choisie \\
    \rowcolor{cfptGrey}
    \texttt{id} & \textbf{unique} pour chacun : c'est le lien avec son \texttt{<label>} \\
    \bottomrule
  \end{tabularx}
  \end{center}
  \vspace{0.25cm}
  \begin{center}
    \small\itshape Deux \texttt{name} différents produisent deux groupes indépendants --- symptôme classique quand « les deux boutons restent allumés ».
  \end{center}
\end{frame}

% ===============================================================
\section{Les attributs qui encadrent la saisie}
% ===============================================================
\partieslide{3. Les attributs qui encadrent la saisie}

\begin{frame}{Le tableau des attributs}
  \vspace{0.1cm}
  \begin{center}
  \scriptsize
  \begin{tabularx}{\textwidth}{lX}
    \toprule
    \rowcolor{accentBlue}
    \hdr{Attribut} & \hdr{Effet} \\
    \midrule
    \rowcolor{cfptGrey}
    \texttt{name} & \textbf{Le nom sous lequel la valeur sera transmise au serveur.} Sans lui, le champ n'est pas envoyé \\
    \texttt{id} & Identifiant unique dans la page ; sert au \texttt{<label>}, au CSS et à JavaScript \\
    \rowcolor{cfptGrey}
    \texttt{value} & Valeur initiale du champ (ou, pour \texttt{radio} et \texttt{checkbox}, la valeur envoyée) \\
    \texttt{placeholder} & Texte d'exemple grisé, qui disparaît dès la première frappe \\
    \rowcolor{cfptGrey}
    \texttt{required} & Le champ doit être rempli pour que le formulaire parte \\
    \texttt{maxlength} / \texttt{minlength} & Bornent le nombre de caractères \\
    \rowcolor{cfptGrey}
    \texttt{min} / \texttt{max} / \texttt{step} & Bornent une valeur numérique ou une date \\
    \texttt{pattern} & Impose un format via une expression régulière \\
    \rowcolor{cfptGrey}
    \texttt{readonly} & Champ visible et \textbf{transmis}, mais non modifiable \\
    \texttt{disabled} & Champ inactif et \textbf{non transmis} au serveur \\
    \bottomrule
  \end{tabularx}
  \end{center}
\end{frame}

\begin{frame}{\texttt{name} et \texttt{id} : deux mondes}
  \begin{center}
  \small
  \begin{tabularx}{\textwidth}{lXX}
    \toprule
    \rowcolor{accentBlue}
    \hdr{} & \hdr{\texttt{id}} & \hdr{\texttt{name}} \\
    \midrule
    \rowcolor{cfptGrey}
    Sert à & désigner l'élément \textbf{dans la page} & nommer la valeur \textbf{pour le serveur} \\
    Utilisé par & le \texttt{<label>}, le CSS, JavaScript & le formulaire au moment de l'envoi \\
    \rowcolor{cfptGrey}
    Unicité & unique dans la page & partagé par les radios d'un groupe \\
    Si absent & impossible de relier un label & \textbf{le champ n'est pas transmis} \\
    \bottomrule
  \end{tabularx}
  \end{center}
  \vspace{0.3cm}
  \begin{alertblock}{La confusion la plus tenace de la semaine}
    Les deux attributs se ressemblent, portent souvent la \textbf{même valeur}, et ne servent pourtant à rien de commun.
  \end{alertblock}
\end{frame}

\begin{frame}{Ce que le serveur recevra en semaine 15}
  \begin{block}{Au moment de l'envoi}
    Le navigateur assemble les champs en une suite de paires \texttt{nom=valeur}. Pour le formulaire de création de sujet :
    \vspace{0.2cm}

    \begin{center}
      \texttt{\small pseudo=alice\&titre=Bonjour\&categorie=general\&message=Salut}
    \end{center}
  \end{block}
  \vspace{0.2cm}
  \begin{exampleblock}{Chaque clé est un attribut \texttt{name}}
    En semaine 15, le serveur Express lira exactement ces clés : \texttt{req.body.pseudo}, \texttt{req.body.titre}\ldots
    \vspace{0.15cm}

    Un champ sans \texttt{name} est \textbf{absent de la liste} : le serveur ne saura même pas qu'il a existé. C'est pour cela qu'on nomme les champs \textbf{dès maintenant}, alors que rien ne les lit encore.
  \end{exampleblock}
\end{frame}

\begin{frame}[fragile]{\texttt{pattern} : imposer un format}
\begin{lstlisting}[language=HTML5]
<label for="pseudo">Pseudo</label>
<input type="text" id="pseudo" name="pseudo"
       pattern="[A-Za-z0-9]{3,15}"
       title="3 a 15 caracteres : lettres et chiffres uniquement"
       required>
\end{lstlisting}
  \vspace{0.2cm}
  \begin{alertblock}{\texttt{pattern} sans \texttt{title} est inutilisable}
    Si la saisie ne correspond pas au motif, le navigateur affiche un message générique du type « Veuillez respecter le format demandé » : l'utilisateur ne sait \textbf{pas quoi corriger}.
    \vspace{0.15cm}

    L'attribut \texttt{title} fournit le texte affiché à la place. Il n'est pas décoratif --- sans lui, le champ devient un piège.
  \end{alertblock}
\end{frame}

% ===============================================================
\section{La validation native et ses limites}
% ===============================================================
\partieslide{4. La validation native et ses limites}

\begin{frame}{Ce que le navigateur fait tout seul}
  Un champ \texttt{required} vide ou un \texttt{type="email"} mal formé \textbf{bloquent l'envoi} et déclenchent une bulle d'avertissement --- sans une ligne de JavaScript.
  \vspace{0.3cm}
  \begin{block}{Ce qu'elle apporte}
    \begin{itemize}
      \item Un retour \textbf{immédiat}, avant tout aller-retour avec le serveur
      \item Des messages traduits dans la langue du navigateur, sans effort
      \item Un premier filtre contre les saisies manifestement erronées
    \end{itemize}
  \end{block}
  \vspace{0.2cm}
  \begin{exampleblock}{Le test à faire sur chaque formulaire du TP}
    Laisse un champ \texttt{required} vide et clique sur le bouton d'envoi. Le navigateur doit refuser, placer le curseur dans le champ fautif et afficher un message.
  \end{exampleblock}
\end{frame}

\begin{frame}{Une validation côté navigateur ne protège rien}
  \begin{alertblock}{C'est une aide à la saisie, pas une sécurité}
    Elle s'exécute sur la machine du \textbf{visiteur}, qui garde la main dessus.
    \vspace{0.15cm}

    \begin{itemize}
      \item Ouvrir les outils de développement et supprimer l'attribut \texttt{required} suffit à envoyer un formulaire vide.
      \item On peut même ne pas utiliser de navigateur du tout : un outil en ligne de commande envoie une requête directement au serveur, sans jamais avoir vu le formulaire.
    \end{itemize}
  \end{alertblock}
  \vspace{0.25cm}
  \begin{block}{Conséquence, valable pour toute l'année}
    Tout contrôle effectué dans le navigateur doit être \textbf{refait sur le serveur}. Mise en \oe uvre en semaine 15 ; en semaine 19, on verra ce qu'un attaquant fait d'un serveur qui a fait confiance au navigateur.
  \end{block}
\end{frame}

\begin{frame}{Les quatre niveaux de contrôle}
  \vspace{0.2cm}
  \begin{center}
  \small
  \begin{tabularx}{\textwidth}{lXl}
    \toprule
    \rowcolor{accentBlue}
    \hdr{Niveau} & \hdr{Ce qu'il assure} & \hdr{Vu en} \\
    \midrule
    \rowcolor{cfptGrey}
    Navigateur (HTML5) & Confort de saisie, erreurs évidentes & semaine 3 \\
    JavaScript (client) & Messages sur mesure, règles combinées & semaine 7 \\
    \rowcolor{cfptGrey}
    Serveur (Express) & \textbf{La seule barrière fiable} & semaine 15 \\
    Base de données & Intégrité des données stockées & semaine 8 \\
    \bottomrule
  \end{tabularx}
  \end{center}
  \vspace{0.4cm}
  \begin{center}
    \small\itshape Les trois premiers niveaux se complètent. Un seul est obligatoire.
  \end{center}
\end{frame}

% ===============================================================
\section{Structurer un formulaire}
% ===============================================================
\partieslide{5. Structurer un formulaire}

\begin{frame}[fragile]{\texttt{<fieldset>} et \texttt{<legend>}}
\begin{lstlisting}[language=HTML5]
<fieldset>
  <legend>Votre message</legend>

  <label for="titre">Titre du sujet</label>
  <input type="text" id="titre" name="titre" maxlength="80" required>

  <label for="message">Message</label>
  <textarea id="message" name="message" rows="6" required></textarea>
</fieldset>
\end{lstlisting}
  \vspace{0.15cm}
  \begin{center}
  \small
  \begin{tabularx}{0.95\textwidth}{lX}
    \toprule
    \rowcolor{accentBlue}
    \hdr{Balise} & \hdr{Rôle} \\
    \midrule
    \rowcolor{cfptGrey}
    \texttt{<fieldset>} & Regroupe des champs qui vont ensemble \\
    \texttt{<legend>} & Le titre du groupe ; \textbf{première balise} à l'intérieur du \texttt{<fieldset>} \\
    \bottomrule
  \end{tabularx}
  \end{center}
\end{frame}

\begin{frame}{Indispensable pour un groupe de boutons radio}
  \begin{alertblock}{Ce qu'entend une personne sans écran}
    Un lecteur d'écran annonce l'étiquette de chaque bouton radio --- « Public », « Réservé aux membres » --- mais \textbf{rien ne dit de quoi il s'agit}.
  \end{alertblock}
  \vspace{0.25cm}
  \begin{exampleblock}{Avec un \texttt{<legend>}}
    La personne entend « \textbf{Visibilité} : Public ». L'information manquante est fournie par le titre du groupe.
  \end{exampleblock}
  \vspace{0.25cm}
  \begin{center}
    \small Pour un champ isolé, le \texttt{<label>} suffit. Pour un \textbf{groupe}, il faut les deux.
  \end{center}
\end{frame}

% ===============================================================
\section{Accessibilité}
% ===============================================================
\partieslide{6. Accessibilité}

\begin{frame}{De quoi parle-t-on ?}
  \begin{block}{Définition --- Accessibilité (a11y)}
    Rendre un site utilisable \textbf{par tous}, y compris les personnes en situation de handicap : malvoyance, cécité, daltonisme, difficultés motrices empêchant l'usage d'une souris, troubles cognitifs.
    \vspace{0.1cm}

    Ce n'est pas une couche que l'on ajoute à la fin : elle se joue à \textbf{80\,\%} dans le choix des balises HTML --- donc maintenant.
  \end{block}
  \vspace{0.1cm}
  \begin{exampleblock}{Trois raisons de s'en préoccuper dès la troisième semaine}
    \begin{itemize}
      \item \textbf{Des gens.} Une page inaccessible exclut des utilisateurs réels, souvent sans que personne s'en aperçoive.
      \item \textbf{La loi.} Les sites publics suisses sont soumis à des exigences d'accessibilité, et la plupart des appels d'offres publics les imposent.
      \item \textbf{La qualité.} Un HTML accessible est un HTML bien structuré : mieux référencé, plus facile à maintenir.
    \end{itemize}
  \end{exampleblock}
\end{frame}

\begin{frame}[fragile]{\texttt{<label>} --- le point le plus important de la semaine}
\begin{lstlisting}[language=HTML5]
<!-- Ecriture explicite : for = id (la plus courante) -->
<label for="pseudo">Votre pseudo</label>
<input type="text" id="pseudo" name="pseudo">

<!-- Ecriture implicite : le label enveloppe le champ -->
<label>
  Votre pseudo
  <input type="text" name="pseudo">
</label>
\end{lstlisting}
  \vspace{0.15cm}
  \begin{alertblock}{Le test d'une seconde}
    \texttt{for} reprend l'\texttt{id}, jamais le \texttt{name}. Pour vérifier : \textbf{clique sur l'étiquette}. Si le curseur se place dans le champ (ou si la case se coche), la liaison fonctionne.
  \end{alertblock}
\end{frame}

\begin{frame}{Un \texttt{placeholder} n'est pas un \texttt{<label>}}
  \begin{alertblock}{Quatre raisons de ne pas supprimer les étiquettes}
    \begin{itemize}
      \item le texte \textbf{disparaît} dès la première frappe --- l'utilisateur ne sait plus ce qu'il remplit ;
      \item son contraste est volontairement faible, donc difficile à lire ;
      \item plusieurs lecteurs d'écran ne l'annoncent pas ;
      \item il n'y a plus de zone cliquable pour atteindre le champ.
    \end{itemize}
  \end{alertblock}
  \vspace{0.3cm}
  \begin{block}{La bonne répartition}
    Le \texttt{placeholder} donne un \textbf{exemple} de saisie (\texttt{ex : alice}).\\
    Le \texttt{<label>} dit \textbf{ce qu'est} le champ.
    \vspace{0.1cm}

    Les deux, pas l'un à la place de l'autre.
  \end{block}
\end{frame}

\begin{frame}{Le texte alternatif dépend du rôle de l'image}
  \vspace{0.15cm}
  \begin{center}
  \small
  \begin{tabularx}{\textwidth}{lX}
    \toprule
    \rowcolor{accentBlue}
    \hdr{Rôle de l'image} & \hdr{Que mettre dans \texttt{alt}} \\
    \midrule
    \rowcolor{cfptGrey}
    Elle porte une information & Ce que l'image \textbf{dit} : \texttt{alt="Trois sujets sans réponse depuis lundi"} \\
    Elle est un lien ou un bouton & La \textbf{destination} ou l'action : \texttt{alt="Retour à l'accueil"} \\
    \rowcolor{cfptGrey}
    Elle est décorative & Un \texttt{alt} \textbf{vide} : \texttt{alt=""}, pour que le lecteur d'écran l'ignore \\
    \bottomrule
  \end{tabularx}
  \end{center}
  \vspace{0.3cm}
  \begin{exampleblock}{Le test du téléphone}
    Imagine que tu décrives la page à quelqu'un au téléphone. Ce que tu dirais de l'image, c'est le \texttt{alt}. Ce que tu passerais sous silence parce que c'est un simple élément de décor, c'est un \texttt{alt=""}.
  \end{exampleblock}
\end{frame}

\begin{frame}{Le contraste des couleurs}
  \begin{block}{Définition --- Rapport de contraste}
    Il compare la luminosité d'un texte et de son fond. Il va de \textbf{1:1} (invisible) à \textbf{21:1} (noir sur blanc).
    \vspace{0.1cm}

    Seuil de référence, niveau \textbf{AA} des règles WCAG : \textbf{4{,}5:1} pour le texte courant, \textbf{3:1} pour les grands titres.
  \end{block}
  \vspace{0.25cm}
  \begin{alertblock}{La couleur ne doit jamais être la seule information}
    Un champ en erreur signalé \textbf{uniquement} par une bordure rouge est invisible pour une personne daltonienne --- environ 8\,\% des hommes.
    \vspace{0.15cm}

    Double systématiquement la couleur d'un second signal : un texte (« Ce champ est obligatoire »), une icône, ou un soulignement. Vaut aussi pour les liens dans un paragraphe.
  \end{alertblock}
\end{frame}

\begin{frame}{La navigation au clavier}
  \vspace{0.1cm}
  \begin{center}
  \small
  \begin{tabularx}{0.95\textwidth}{lX}
    \toprule
    \rowcolor{accentBlue}
    \hdr{Touche} & \hdr{Action} \\
    \midrule
    \rowcolor{cfptGrey}
    \texttt{Tab} & Passer à l'élément interactif suivant \\
    \texttt{Maj + Tab} & Revenir au précédent \\
    \rowcolor{cfptGrey}
    \texttt{Entrée} & Activer un lien ou envoyer un formulaire \\
    \texttt{Espace} & Cocher une case, activer un bouton \\
    \rowcolor{cfptGrey}
    Flèches & Se déplacer dans un groupe de radios ou une liste déroulante \\
    \bottomrule
  \end{tabularx}
  \end{center}
  \vspace{0.25cm}
  \begin{exampleblock}{Le test des trente secondes}
    Clique dans la barre d'adresse, puis parcours toute la page avec \texttt{Tab} sans jamais toucher la souris. Trois choses doivent être vraies : tu \textbf{vois toujours} où tu te trouves ; l'ordre suit la \textbf{lecture naturelle} ; tu peux atteindre \textbf{tous} les liens, champs et boutons.
  \end{exampleblock}
\end{frame}

\begin{frame}[fragile]{Deux règles à retenir pour la suite}
  \begin{block}{L'ordre de tabulation suit l'ordre du \textbf{code}}
    Pas l'ordre visuel. Si le CSS déplace un élément (semaines 4 et 5), la navigation clavier continuera de suivre le HTML.
    \vspace{0.1cm}

    Une raison de plus d'écrire le HTML dans un ordre logique dès maintenant.
  \end{block}
  \vspace{0.2cm}
\begin{lstlisting}[language=CSS]
/* A ne jamais ecrire */
button:focus { outline: none; }
\end{lstlisting}
  \vspace{0.1cm}
  \begin{alertblock}{Ne supprime jamais le cadre de focus}
    On trouve souvent cette ligne dans du CSS repris en ligne : à elle seule, elle rend un site inutilisable au clavier. Si le cadre par défaut ne plaît pas, on le \textbf{remplace} par un autre style visible --- on ne l'efface pas.
  \end{alertblock}
\end{frame}

\begin{frame}{Vérifier son travail}
  \vspace{0.1cm}
  \begin{center}
  \small
  \begin{tabularx}{\textwidth}{lX}
    \toprule
    \rowcolor{accentBlue}
    \hdr{Outil} & \hdr{Ce qu'il apporte} \\
    \midrule
    \rowcolor{cfptGrey}
    Onglet \textit{Accessibilité} des DevTools & L'arbre tel que le perçoit un lecteur d'écran, et les contrastes insuffisants \\
    Lighthouse (Chrome) & Un rapport noté, avec la liste des problèmes détectés \\
    \rowcolor{cfptGrey}
    \texttt{validator.w3.org} & Les \texttt{alt} manquants, les \texttt{for} orphelins, les \texttt{id} en double \\
    Le clavier & Le test le plus rapide, et celui qu'aucun outil ne remplace \\
    \bottomrule
  \end{tabularx}
  \end{center}
  \vspace{0.3cm}
  \begin{alertblock}{Aucun outil automatique ne suffit}
    Un rapport Lighthouse à 100 ne signifie pas qu'une page est accessible : une machine sait vérifier qu'un \texttt{alt} \textbf{existe}, pas qu'il \textbf{décrit correctement} l'image. Ces outils repèrent les erreurs mécaniques ; le reste relève du jugement.
  \end{alertblock}
\end{frame}

\begin{frame}{Les huit réflexes de la semaine}
  \begin{columns}[T]
    \begin{column}{0.49\textwidth}
      \begin{block}{Formulaire}
        \begin{itemize}
          \item Chaque \texttt{<input>} a un \texttt{<label>} relié par \texttt{for} / \texttt{id}
          \item Chaque champ à transmettre a un \texttt{name}
          \item Les groupes de radios sont dans un \texttt{<fieldset>} avec \texttt{<legend>}
          \item Le \texttt{placeholder} complète le label, il ne le remplace pas
        \end{itemize}
      \end{block}
    \end{column}
    \begin{column}{0.49\textwidth}
      \begin{exampleblock}{Page}
        \begin{itemize}
          \item Chaque \texttt{<img>} a un \texttt{alt} adapté à son rôle
          \item L'information n'est jamais portée par la seule couleur
          \item La page se parcourt entièrement au clavier, focus visible
          \item La page passe le validateur W3C sans erreur
        \end{itemize}
      \end{exampleblock}
    \end{column}
  \end{columns}
\end{frame}

% ===============================================================
% Clôture
% ===============================================================
\begin{frame}{Résumé de la semaine}
  \begin{columns}[T]
    \begin{column}{0.5\textwidth}
      \begin{block}{Ce qu'on a vu}
        \begin{itemize}
          \item \texttt{<form>}, \texttt{<input>}, \texttt{<textarea>}, \texttt{<select>}, \texttt{<button>}
          \item Le \texttt{type} change tout : contrôle et clavier mobile
          \item \texttt{name} pour le serveur, \texttt{id} pour la page
          \item La validation native n'est pas une sécurité
          \item \texttt{<label>}, \texttt{<fieldset>}, \texttt{<legend>}
          \item \texttt{alt}, contraste, navigation au clavier
        \end{itemize}
      \end{block}
    \end{column}
    \begin{column}{0.5\textwidth}
      \begin{alertblock}{À faire pour la semaine 4}
        \begin{itemize}
          \item Les \textbf{quatre formulaires} du MiniForum terminés (fiche TP)
          \item Chaque page passe \texttt{validator.w3.org} sans erreur
          \item Le site est entièrement parcourable au clavier
          \item Le tout poussé sur GitHub, en plusieurs commits
          \item Faire le quiz Moodle
        \end{itemize}
      \end{alertblock}
    \end{column}
  \end{columns}
\end{frame}

\begin{frame}{Fil rouge --- MiniForum : tous les formulaires en place}
  \begin{block}{Où en est le projet}
    Le MiniForum contient désormais l'intégralité de ses formulaires : création de sujet, réponse à un sujet, inscription et connexion.
  \end{block}
  \vspace{0.25cm}
  \begin{exampleblock}{Ce qui ne bougera plus}
    Ils ne font encore rien --- mais leur structure ne changera plus jusqu'à la fin du cours.
    \vspace{0.15cm}

    Les \texttt{name} choisis cette semaine sont ceux que le serveur Express lira en \textbf{semaine 15}. Les \texttt{id} sont ceux que JavaScript utilisera dès la \textbf{semaine 7}.
  \end{exampleblock}
  \vspace{0.2cm}
  \begin{center}
    \small Trois semaines de HTML sont terminées. La semaine 4 ouvre le \textbf{CSS} : la structure va enfin recevoir sa mise en forme.
  \end{center}
\end{frame}

{
\setbeamertemplate{footline}{}
\begin{frame}[plain]
  \begin{tikzpicture}[remember picture, overlay]
    \fill[accentBlue] (current page.north west) rectangle (current page.south east);
    \node[text=white, font=\Huge\bfseries] at ([yshift=0.6cm]current page.center) {Des questions ?};
    \fill[cfptYellow] ([yshift=-0.2cm]current page.center -| current page.west)
      rectangle ([yshift=-0.32cm]current page.center -| current page.east);
    \node[text=lightBlue, font=\large] at ([yshift=-1.3cm]current page.center)
      {Prochaine étape : la fiche TP --- les quatre formulaires du MiniForum};
  \end{tikzpicture}
\end{frame}
}

\end{document}
